% GENERATED FILE. Edit canonical lessons, then run npm run generate:books.
% canonical-source-hash: fnv1a64:b94cd50b1678c58e
% canonical-lessons: PA-C105-kattak, PA-C105-magghar, PA-C105-poh, PA-C105-magh, PA-C105-phaggan

\chapter{Kattak, Magghar, Poh, Magh, Phaggan}
\label{ch:pa-kattak-magghar-poh-magh-phaggan}
\hlchaptermodality{105}

\begin{chapteropening}
\textbf{By the end of this chapter:} \emph{I can say Kattak, Magghar, Poh, Magh and Phaggan.}

Complete the last lesson of chapter 105: I can say Kattak, Magghar, Poh, Magh and Phaggan.
\end{chapteropening}

\section[\texorpdfstring{\pa{ਕੱਤਕ} (kattak)}{ਕੱਤਕ (kattak)}]{\pa{ਕੱਤਕ} (kattak) — Kattak, a Punjabi month}

\label{lesson:PA-C105-kattak}



\begin{quote}
Before the new one: say the Punjabi for Bhadon, then the Punjabi for Assu.
\end{quote}

\subsection*{You'll want to know: \pa{ਕੱਤਕ}}
\textbf{\pa{ਕੱਤਕ}} — \emph{kattak} — \textquotedblleft{}Kattak, a Punjabi month\textquotedblright{}.

\begin{cousinweb}
\textbf{\pa{ਕੱਤਕ}} (\emph{kattak}) goes back to Sanskrit \emph{kārttika}, \textquotedblleft{}the month October-November\textquotedblright{}. CDIAL gives no Punjabi form, but lists Lahnda \emph{kattẽ} under it (CDIAL 3070).
\end{cousinweb}

\begin{grammarlens}[title={What you've built}]
One more word for this chapter.
\end{grammarlens}

\subsection*{Your turn}
\begin{itemize}
\raggedright
  \item \emph{Say it:} \emph{kattak}
  \item \emph{Say it:} \emph{kattak}, once more
  \item \emph{Say it:} say \emph{kattak}
  \item \emph{Recall:} say the Punjabi for Bhadon, then the Punjabi for Assu, then say \emph{kattak} again
\end{itemize}

\begin{tcolorbox}[breakable,colback=teal!4,colframe=teal!35!black,title={Before you move on}]
What does \pa{ਕੱਤਕ} mean? (\textquotedblleft{}Kattak, a Punjabi month\textquotedblright{}.) Say it once more.
\end{tcolorbox}
\section[\texorpdfstring{\pa{ਮੱਘਰ} (magghar)}{ਮੱਘਰ (magghar)}]{\pa{ਮੱਘਰ} (magghar) — Magghar, a Punjabi month}

\label{lesson:PA-C105-magghar}



\begin{quote}
Before the new one: say the Punjabi for Assu, then the Punjabi for Kattak.
\end{quote}

\subsection*{You'll want to know: \pa{ਮੱਘਰ}}
\textbf{\pa{ਮੱਘਰ}} — \emph{magghar} — \textquotedblleft{}Magghar, a Punjabi month\textquotedblright{}.

\begin{cousinweb}
\textbf{\pa{ਮੱਘਰ}} (\emph{magghar}) is a tadbhava, a word inherited from Sanskrit and reshaped by sound change. It comes from \emph{mārgaśiras}, \textquotedblleft{}the month November-December\textquotedblright{} (CDIAL 10076).
\end{cousinweb}

\begin{grammarlens}[title={What you've built}]
One more word for this chapter.
\end{grammarlens}

\subsection*{Your turn}
\begin{itemize}
\raggedright
  \item \emph{Say it:} \emph{magghar}
  \item \emph{Say it:} \emph{magghar}, once more
  \item \emph{Say it:} say \emph{magghar}
  \item \emph{Recall:} say the Punjabi for Assu, then the Punjabi for Kattak, then say \emph{magghar} again
\end{itemize}

\begin{tcolorbox}[breakable,colback=teal!4,colframe=teal!35!black,title={Before you move on}]
What does \pa{ਮੱਘਰ} mean? (\textquotedblleft{}Magghar, a Punjabi month\textquotedblright{}.) Say it once more.
\end{tcolorbox}
\section[\texorpdfstring{\pa{ਪੋਹ} (poh)}{ਪੋਹ (poh)}]{\pa{ਪੋਹ} (poh) — Poh, a Punjabi month}

\label{lesson:PA-C105-poh}



\begin{quote}
Before the new one: say the Punjabi for Kattak, then the Punjabi for Magghar.
\end{quote}

\subsection*{You'll want to know: \pa{ਪੋਹ}}
\textbf{\pa{ਪੋਹ}} — \emph{poh} — \textquotedblleft{}Poh, a Punjabi month\textquotedblright{}.

\begin{cousinweb}
\textbf{\pa{ਪੋਹ}} (\emph{poh}) is a tadbhava, a word inherited from Sanskrit and reshaped by sound change. It comes from \emph{pauṣa}, \textquotedblleft{}the month December-January\textquotedblright{} (CDIAL 8424).
\end{cousinweb}

\begin{grammarlens}[title={What you've built}]
One more word for this chapter.
\end{grammarlens}

\subsection*{Your turn}
\begin{itemize}
\raggedright
  \item \emph{Say it:} \emph{poh}
  \item \emph{Say it:} \emph{poh}, once more
  \item \emph{Say it:} say \emph{poh}
  \item \emph{Recall:} say the Punjabi for Kattak, then the Punjabi for Magghar, then say \emph{poh} again
\end{itemize}

\begin{tcolorbox}[breakable,colback=teal!4,colframe=teal!35!black,title={Before you move on}]
What does \pa{ਪੋਹ} mean? (\textquotedblleft{}Poh, a Punjabi month\textquotedblright{}.) Say it once more.
\end{tcolorbox}
\section[\texorpdfstring{\pa{ਮਾਘ} (māgh)}{ਮਾਘ (māgh)}]{\pa{ਮਾਘ} (māgh) — Magh, a Punjabi month}

\label{lesson:PA-C105-magh}



\begin{quote}
Before the new one: say the Punjabi for Magghar, then the Punjabi for Poh.
\end{quote}

\subsection*{You'll want to know: \pa{ਮਾਘ}}
\textbf{\pa{ਮਾਘ}} — \emph{māgh} — \textquotedblleft{}Magh, a Punjabi month\textquotedblright{}.

\begin{cousinweb}
\textbf{\pa{ਮਾਘ}} (\emph{māgh}) is a semi-tatsama: Sanskrit \emph{māgha}, \textquotedblleft{}the month January-February\textquotedblright{}, borrowed after the old sound changes had run their course and then lightly reshaped (CDIAL 9993). CDIAL lists the inherited Punjabi form as \emph{māh}.
\end{cousinweb}

\begin{grammarlens}[title={What you've built}]
One more word for this chapter.
\end{grammarlens}

\subsection*{Your turn}
\begin{itemize}
\raggedright
  \item \emph{Say it:} \emph{māgh}
  \item \emph{Say it:} \emph{māgh}, once more
  \item \emph{Say it:} say \emph{māgh}
  \item \emph{Recall:} say the Punjabi for Magghar, then the Punjabi for Poh, then say \emph{māgh} again
\end{itemize}

\begin{tcolorbox}[breakable,colback=teal!4,colframe=teal!35!black,title={Before you move on}]
What does \pa{ਮਾਘ} mean? (\textquotedblleft{}Magh, a Punjabi month\textquotedblright{}.) Say it once more.
\end{tcolorbox}
\section[\texorpdfstring{\pa{ਫੱਗਣ} (phaggaṇ)}{ਫੱਗਣ (phaggaṇ)}]{\pa{ਫੱਗਣ} (phaggaṇ) — Phaggan, a Punjabi month}

\label{lesson:PA-C105-phaggan}



\begin{quote}
Before the new one: say the Punjabi for Poh, then the Punjabi for Magh.
\end{quote}

\subsection*{You'll want to know: \pa{ਫੱਗਣ}}
\textbf{\pa{ਫੱਗਣ}} — \emph{phaggaṇ} — \textquotedblleft{}Phaggan, a Punjabi month\textquotedblright{}.

\begin{cousinweb}
\textbf{\pa{ਫੱਗਣ}} (\emph{phaggaṇ}) is a tadbhava, a word inherited from Sanskrit and reshaped by sound change. It comes from \emph{phalguna}, \textquotedblleft{}reddish; the month February-March\textquotedblright{} (CDIAL 9065).
\end{cousinweb}

\begin{grammarlens}[title={What you've built}]
One more word for this chapter.
\end{grammarlens}

\subsection*{Your turn}
\begin{itemize}
\raggedright
  \item \emph{Say it:} \emph{phaggaṇ}
  \item \emph{Say it:} \emph{phaggaṇ}, once more
  \item \emph{Say it:} say \emph{phaggaṇ}
  \item \emph{Recall:} say the Punjabi for Poh, then the Punjabi for Magh, then say \emph{phaggaṇ} again
\end{itemize}

\begin{tcolorbox}[breakable,colback=teal!4,colframe=teal!35!black,title={Before you move on}]
What does \pa{ਫੱਗਣ} mean? (\textquotedblleft{}Phaggan, a Punjabi month\textquotedblright{}.) Say it once more.
\end{tcolorbox}
